\chapter{Sorting and greedy algorithms}
\label{ch:greedy}
\index{greedy algorithm}

Half the problems at this level are won by two moves you already own: put
the input in a useful order, then take whatever looks best next and never
look back. The first move is \texttt{std::sort}; the second is a
\emph{greedy} algorithm. Sorting is safe --- it always produces a sorted
array. Greedy is \emph{not} safe: it is a claim about the problem, and the
claim needs a proof. This chapter teaches the tools (comparators,
coordinate compression, \texttt{set} as a ``best available'' structure),
the two standard proofs (the exchange argument and ``greedy stays
ahead''), and the discipline of hunting counter-examples before trusting
a greedy rule. After it, you will sort fluently, compress fluently, and
--- more importantly --- know which greedy ideas deserve your trust.

\section{Sorting as a tool}
\label{sec:sorting}
\index{sorting!std::sort@\texttt{std::sort}}

\texttt{std::sort} runs in $O(n\log n)$ worst case (introsort: quicksort
with a heapsort fallback) and is \emph{unstable}: equal elements may
reorder. When equal-order matters, \texttt{stable\_sort} preserves it ---
a detail that costs a contest point exactly once, the time you forget it.
A comparator must define a \emph{strict weak ordering}: irreflexive,
asymmetric, transitive. The classic violation is \texttt{<=} instead of
\texttt{<}: equal elements then compare ``less'' than each other, the
invariant breaks, and \texttt{sort} can read out of bounds --- real
undefined behaviour, not just a wrong answer\index{comparator}.

\begin{lstlisting}[label={lst:sortcmp},
caption={\texttt{code/sort\_cmp.cpp}: lambda comparator plus structured
bindings. The comparator asks one question --- ``does $a$ come before
$b$?'' --- never ``are they equal?''}]
\end{lstlisting}
\lstinputlisting{code/sort_cmp.cpp}
Its real output --- note the tie broken by name:
\begin{lstlisting}[numbers=none]
bella 9
cato 9
anni 7
dele 4
\end{lstlisting}

The second tool is \emph{coordinate compression}\index{coordinate
compression}: when only the \emph{order} of values matters, replace each
by its rank among the distinct values. Values up to $4\cdot10^9$ become
ranks $0\dots k{-}1$ that index a Fenwick tree or a small array ---
essential from Chapter~12 on.

\begin{lstlisting}[label={lst:compress},
caption={\texttt{code/compress\_demo.cpp} using
\texttt{code/lib/compress.h}. Ranks preserve every comparison
($<$, $=$, $>$); nothing else survives, by design.}]
\end{lstlisting}
\lstinputlisting{code/compress_demo.cpp}
\begin{lstlisting}[numbers=none]
4000000000 -> 2
10 -> 0
10 -> 0
80 -> 1
4000000000 -> 2
distinct values: 3
\end{lstlisting}
Figure~\ref{fig:compress} draws the same run.

\begin{figure}[ht]\centering
% fig13-compress.tex -- Figure 3.4: values -> ranks for {4e9,10,10,80,4e9}.
% fstyle.tex -- shared TikZ styles for the figures of this book (palette from lifeos.sty).
\tikzset{%
  cell/.style={draw=ink, minimum width=7mm, minimum height=6.5mm, inner sep=0pt, font=\footnotesize, fill=white},
  cellh/.style={cell, fill=accentsoft, draw=accent},
  cella/.style={cell, fill=amberbg, draw=amber},
  cellg/.style={cell, fill=goodbg, draw=good},
  cellx/.style={cell, fill=graybg, draw=slate, text=slate},
  vtx/.style={circle, draw=ink, minimum size=6.5mm, inner sep=0pt, font=\footnotesize, fill=white},
  vtxh/.style={vtx, draw=accent, fill=accentsoft},
  vtxa/.style={vtx, draw=amber, fill=amberbg},
  vtxg/.style={vtx, draw=good, fill=goodbg},
  vtxd/.style={vtx, draw=slate, fill=graybg, text=slate},
  lab/.style={font=\footnotesize, text=slate},
  labA/.style={font=\footnotesize, text=accent},
  labM/.style={font=\footnotesize, text=amber},
  idx/.style={font=\tiny, text=slate},
  arr/.style={-{stealth}, thick, ink},
  arrA/.style={-{stealth}, thick, accent},
  arrM/.style={-{stealth}, thick, amber},
  arrG/.style={-{stealth}, thick, good},
  edge/.style={thick, ink},
  edgeA/.style={thick, accent},
  edgeM/.style={thick, amber},
}%
% \cellrow{prefix}{x0}{y}{style}{v0,v1,...}: a row of cells named prefix0, prefix1, ... (pitch 0.75cm)
\providecommand{\cellrow}[5]{\foreach \v [count=\i from 0] in {#5} {\node[#4] (#1\i) at (#2+\i*0.75,#3) {\v};}}%
% \cellidx{prefix}{n}{start}: index labels under cells prefix0 .. prefix(n-1)
\providecommand{\cellidx}[3]{\foreach \i in {0,...,\numexpr#2-1\relax} {\node[idx, below=0pt of #1\i] {\the\numexpr\i+#3\relax};}}%
% \fitfig{<tikzpicture>}: shrink the picture to \linewidth only when it would overflow.
\ifdefined\fitbox\else\newsavebox{\fitbox}\fi%
\providecommand{\fitfig}[1]{\sbox{\fitbox}{#1}\ifdim\wd\fitbox>\linewidth\resizebox{\linewidth}{!}{\usebox{\fitbox}}\else\usebox{\fitbox}\fi}%

\begin{tikzpicture}[font=\footnotesize]
\node[lab, anchor=east] at (-0.2,1.1) {values};
\node[lab, anchor=east] at (-0.2,-0.3) {ranks};
\foreach \i/\v in {0/4000000000,1/10,2/10,3/80,4/4000000000}
  \node[cellx, minimum width=2.2cm] (v\i) at (\i*2.4+1,1.1) {\v};
\foreach \i/\r in {0/2,1/0,2/0,3/1,4/2}
  \node[cellh, minimum width=0.9cm] (r\i) at (\i*2.4+1,-0.3) {\r};
\foreach \i in {0,...,4} \draw[arrA] (v\i.south) -- (r\i.north);
\node[draw=ink, fill=white, inner sep=4pt] at (6.4,-1.45) {sorted distinct values: $[10,\;80,\;4{\cdot}10^9]$};
\node[lab, align=center] at (6.4,-2.4)
  {rank $=$ position among the sorted distinct values; $<$, $=$, $>$ survive, magnitudes do not};
\end{tikzpicture}

\caption{Values $\to$ ranks. Notice what is thrown away: magnitudes,
gaps, units. If your algorithm ever reads a difference of two values,
compression is the wrong tool.}
\label{fig:compress}
\end{figure}

\begin{exercise}[easy]
Sort $\{(3,\text{a}),(1,\text{c}),(3,\text{b})\}$ by first element only.
Which of the two possible outputs can \texttt{sort} produce, and which
does \texttt{stable\_sort} guarantee?
\end{exercise}

\section{The greedy method}
\label{sec:greedy}
\index{greedy algorithm!definition}

\begin{definition}[Greedy algorithm]
A \emph{greedy algorithm} builds a solution in steps, at each step taking
the locally best choice under some fixed rule --- and never revisiting
the choice. Correctness is not part of the definition.
\end{definition}

Greedy is fast to write and fast to run --- that is the appeal. The cost
is outsourced to a proof: ``locally best is globally right'' is a theorem
about the problem, not about your code. Chapter~2's complete search was
the pessimist's answer (check everything); greedy is the optimist's
answer (commit at every step). The skill is knowing which optimism is
justified.

\begin{example}[CSES \emph{Ferris Wheel}]
Each gondola carries at most two children of total weight $\le x$.
Minimise gondolas. \emph{Strategy.} The heaviest child must ride
somewhere; the only question is whether anyone can join. If the lightest
does not fit with them, nobody does --- so pair heaviest with lightest
whenever possible, else heaviest rides alone.
\end{example}
\lstinputlisting{code/ferris_wheel.cpp}
Real output on the CSES sample (\texttt{4 10 / 7 2 3 9}):
\begin{lstlisting}[numbers=none]
3
\end{lstlisting}
Why does pairing heaviest with lightest work? If the lightest fits,
using any heavier partner instead cannot do better --- the lightest is
\emph{at least} as hard to place as whoever replaces it. That is the
``stays ahead'' instinct of Section~\ref{sec:proofs}.

\section{How to argue correctness}
\label{sec:proofs}
\index{exchange argument}\index{greedy stays ahead}

Two proofs cover most contest greedies:

\begin{itemize}
\item \textbf{Exchange argument.} Take an optimal solution $O$ and the
greedy solution $G$. Find the first place they differ, and show $O$ can
be \emph{exchanged} into $G$'s choice without making things worse.
Repeating morphs $O$ into something no worse that ends with $G$'s
choices --- so greedy is optimal too.
\item \textbf{Greedy stays ahead.} Maintain an invariant like ``after $k$
steps, greedy's partial solution is at least as good as any other partial
solution'', proved by induction on $k$.
\end{itemize}

\begin{example}[CSES \emph{Movie Festival} --- the exchange argument]
\label{ex:movie}
Choose a maximum set of pairwise non-overlapping movies
$(a_i, b_i)$. \emph{Strategy.} Sort by \emph{ending} time and always take
the movie that frees you earliest. Figure~\ref{fig:exchange} is the proof
in one picture.
\end{example}
\lstinputlisting{code/movie_festival.cpp}
Real output on the CSES sample (\texttt{3 5 / 4 9 / 5 8}):
\begin{lstlisting}[numbers=none]
2
\end{lstlisting}

\begin{theorem}[End-time greedy is optimal]
\label{thm:movie}
Let $G = (g_1, \dots, g_t)$ be the schedule produced by always taking the
compatible movie with the smallest end time, and $O = (o_1, \dots, o_m)$
any optimal schedule (both sorted by end time). Then $t = m$.
\end{theorem}
\begin{proof}
Let $j$ be the first index where $G$ and $O$ differ (if $t \ge m$ and they
agree on $m$ entries, we are done --- greedy already has $m$). Since
$g_j$ is the compatible movie with the smallest end, $b(g_j) \le b(o_j)$.
Both fit after the common prefix: $g_j$ and $o_j$ each start no earlier
than $o_{j-1}$'s end. Replace $o_j$ by $g_j$ in $O$: the new schedule is
still feasible --- $o_{j+1}$ starts after $b(o_j) \ge b(g_j)$, so it also
clears $g_j$ --- and has the same size $m$. Repeating the exchange turns
$O$ into a schedule that starts with $G$'s first $t$ movies and still has
$m$ entries. If $m > t$, the tail movie after $g_t$ was still available
when greedy stopped --- contradiction. So $m = t$.
\end{proof}

The picture worth keeping: every exchange \emph{frees time}, never
consumes it --- Figure~\ref{fig:exchange}. And the wrong greedy rule
fails concretely: on $\{(0,10),(1,2),(3,4)\}$, earliest-\emph{start}
takes the long movie and stops at $1$, while the end rule takes both
short ones --- Listing~\ref{lst:sortcmp}'s program prints $1$ and $2$
respectively, and Figure~\ref{fig:timeline} draws both runs.

\begin{figure}[ht]\centering
% fig10-interval.tex -- Figure 3.1: interval scheduling, same instance under two greedy orders.
% Counterexample instance: (0,10) (1,2) (3,4).
\begin{tikzpicture}[x=0.78cm,y=0.62cm,font=\footnotesize,
  take/.style={draw=accent,fill=accentsoft,rounded corners=1pt,minimum height=0.5cm,inner sep=0pt},
  skip/.style={draw=slate,fill=graybg,dashed,rounded corners=1pt,minimum height=0.5cm,inner sep=0pt}]
\draw[->,ink] (0,0) -- (10.8,0) node[below left]{time};
\foreach \t in {0,1,...,10} {\draw[ink] (\t,-0.08) -- (\t,0.08); \node[below,font=\tiny,text=slate] at (\t,-0.1){\t};}
% by start time
\node[anchor=east,align=right] at (-0.2,4.0) {by earliest start\\[-1pt]{\scriptsize picks (0,10)}};
\node[take,minimum width=7.8cm,anchor=west] at (0,4.0) {$(0,10)$};
\node[skip,minimum width=0.78cm,anchor=west] at (1,3.15) {\scriptsize$(1,2)$};
\node[skip,minimum width=0.78cm,anchor=west] at (3,3.15) {\scriptsize$(3,4)$};
\node[amber,anchor=west] at (10.2,4.0) {$=1$ job};
% by end time
\node[anchor=east,align=right] at (-0.2,1.75) {by earliest end\\[-1pt]{\scriptsize picks (1,2),(3,4)}};
\node[take,minimum width=0.78cm,anchor=west] at (1,1.75) {\scriptsize$(1,2)$};
\node[take,minimum width=0.78cm,anchor=west] at (3,1.75) {\scriptsize$(3,4)$};
\node[skip,minimum width=7.8cm,anchor=west] at (0,0.95) {\scriptsize$(0,10)$ overlaps both};
\node[accent,anchor=west] at (10.2,1.75) {$=2$ jobs};
\end{tikzpicture}

\caption{Same instance, two greedy orders. Notice the mechanism of the
failure: the start rule burns the whole timeline on its first pick; the
end rule's picks are exactly the ones that leave room.}
\label{fig:timeline}
\end{figure}

\begin{figure}[ht]\centering
\begin{tikzpicture}[x=0.62cm,y=0.62cm]
% F3.2: the exchange — first index where optimal O and greedy G differ.
\begin{scope}[font=\footnotesize]
\node[anchor=east] at (-0.2,1.0) {optimal $O$:};
\node[anchor=east] at (-0.2,2.6) {greedy $G$:};
% shared prefix
\draw[fill=graybg,draw=slate] (0,0.75) rectangle (1.4,1.25);
\draw[fill=graybg,draw=slate] (0,2.35) rectangle (1.4,2.85);
\node at (1.9,1.0) {$\cdots$}; \node at (1.9,2.6) {$\cdots$};
% first differing pair: g_j ends no later than o_j
\draw[fill=accentsoft,draw=accent,thick] (2.4,2.35) rectangle (4.4,2.85)
  node[midway,font=\tiny]{$g_j$};
\draw[draw=slate,thick] (2.4,0.75) rectangle (5.6,1.25)
  node[midway,font=\tiny]{$o_j$};
\draw[fill=graybg,draw=slate] (6.0,0.75) rectangle (7.6,1.25);
\draw[fill=graybg,draw=slate] (4.8,2.35) rectangle (6.4,2.85);
\draw[{stealth}-{stealth},amber,thick] (4.4,-0.15) -- (5.6,-0.15)
  node[midway,below,font=\footnotesize,text=amber]{$g_j$ frees this much earlier};
\draw[{stealth}->,accent,thick] (3.4,2.3) -- (3.4,1.35);
\end{scope}
\end{tikzpicture}

\caption{The exchange step. At the first difference, greedy's pick ends
no later; swapping it in preserves feasibility \emph{and} size. ``No
later'' is the whole lemma --- everything else is bookkeeping.}
\label{fig:exchange}
\end{figure}

\begin{exercise}[medium]
In the exchange above, where exactly is
``$b(g_j) \le b(o_j)$'' used --- and which step breaks if the greedy rule
picked by earliest \emph{start} instead?
\end{exercise}

\begin{example}[CSES \emph{Stick Lengths} --- the pairing proof]
\label{ex:sticks}
Make all $p_i$ equal at cost $\sum |p_i - m|$. Which $m$?
\emph{Strategy.} Sort and pair the extremes. If $m$ lies between a
smallest and a largest element, that pair contributes exactly their own
gap regardless of $m$'s position inside it --- so the middle element
costs nothing to aim at, and the median is optimal.
\begin{proof}[Proof sketch]
Sort $p$. For each pair $(p_i, p_{n+1-i})$ with $i \le n/2$: if
$p_i \le m \le p_{n+1-i}$, the pair contributes
$(m - p_i) + (p_{n+1-i} - m) = p_{n+1-i} - p_i$, a constant. Moving $m$
outside the pair raises the sum by twice the overshoot. Every paired $m$
between the two middle elements is optimal; the median
$p_{\lceil n/2 \rceil}$ is one.
\end{proof}
\end{example}
\lstinputlisting{code/stick_lengths.cpp}
Real output on the CSES sample (\texttt{2 3 1 5 2}):
\begin{lstlisting}[numbers=none]
5
\end{lstlisting}
Figure~\ref{fig:median} shows the pairing on this data; the stress
harness in \texttt{code/test\_ch03.sh} checked the claim against an
exhaustive minimiser on 300 random arrays --- median won every round.

\begin{figure}[ht]\centering
% fig12-median.tex -- Figure 3.3: median pairing on the CSES sample {1,2,2,3,5}, m=p_3=2.
% fstyle.tex -- shared TikZ styles for the figures of this book (palette from lifeos.sty).
\tikzset{%
  cell/.style={draw=ink, minimum width=7mm, minimum height=6.5mm, inner sep=0pt, font=\footnotesize, fill=white},
  cellh/.style={cell, fill=accentsoft, draw=accent},
  cella/.style={cell, fill=amberbg, draw=amber},
  cellg/.style={cell, fill=goodbg, draw=good},
  cellx/.style={cell, fill=graybg, draw=slate, text=slate},
  vtx/.style={circle, draw=ink, minimum size=6.5mm, inner sep=0pt, font=\footnotesize, fill=white},
  vtxh/.style={vtx, draw=accent, fill=accentsoft},
  vtxa/.style={vtx, draw=amber, fill=amberbg},
  vtxg/.style={vtx, draw=good, fill=goodbg},
  vtxd/.style={vtx, draw=slate, fill=graybg, text=slate},
  lab/.style={font=\footnotesize, text=slate},
  labA/.style={font=\footnotesize, text=accent},
  labM/.style={font=\footnotesize, text=amber},
  idx/.style={font=\tiny, text=slate},
  arr/.style={-{stealth}, thick, ink},
  arrA/.style={-{stealth}, thick, accent},
  arrM/.style={-{stealth}, thick, amber},
  arrG/.style={-{stealth}, thick, good},
  edge/.style={thick, ink},
  edgeA/.style={thick, accent},
  edgeM/.style={thick, amber},
}%
% \cellrow{prefix}{x0}{y}{style}{v0,v1,...}: a row of cells named prefix0, prefix1, ... (pitch 0.75cm)
\providecommand{\cellrow}[5]{\foreach \v [count=\i from 0] in {#5} {\node[#4] (#1\i) at (#2+\i*0.75,#3) {\v};}}%
% \cellidx{prefix}{n}{start}: index labels under cells prefix0 .. prefix(n-1)
\providecommand{\cellidx}[3]{\foreach \i in {0,...,\numexpr#2-1\relax} {\node[idx, below=0pt of #1\i] {\the\numexpr\i+#3\relax};}}%
% \fitfig{<tikzpicture>}: shrink the picture to \linewidth only when it would overflow.
\ifdefined\fitbox\else\newsavebox{\fitbox}\fi%
\providecommand{\fitfig}[1]{\sbox{\fitbox}{#1}\ifdim\wd\fitbox>\linewidth\resizebox{\linewidth}{!}{\usebox{\fitbox}}\else\usebox{\fitbox}\fi}%

\begin{tikzpicture}[x=1.15cm,y=1cm,font=\footnotesize]
\draw[->,ink] (0.4,0) -- (6.1,0) node[right,font=\scriptsize]{$x$};
\foreach \t in {1,...,6} {\draw[ink] (\t,-0.07) -- (\t,0.07); \node[below,idx] at (\t,-0.08) {\t};}
\foreach \x/\c in {1/ink,2/ink,3/ink,5/ink} \fill[\c] (\x,0) circle (2.2pt);
\fill[accent] (2,0.22) circle (2.2pt);
\node[labA, anchor=east] at (1.8,0.3) {$p_2,p_3{=}m$};
% outer pair (p1,p5)
\draw[amber, thick] (1,1.0) -- (5,1.0);
\draw[amber, thick] (1,0.9) -- (1,1.1); \draw[amber, thick] (5,0.9) -- (5,1.1);
\node[labM, above=1pt] at (3,1.0) {$(p_1,p_5)=(1,5)$: $|1-m|+|5-m|=4$};
% inner pair (p2,p4)
\draw[good, thick] (2,0.62) -- (3,0.62);
\draw[good, thick] (2,0.52) -- (2,0.72); \draw[good, thick] (3,0.52) -- (3,0.72);
\node[font=\scriptsize, text=good, anchor=west] at (3.1,0.62) {$(p_2,p_4)=(2,3)$: cost $1$};
\node[lab, text width=9.5cm, align=center] at (3.5,-1.0)
  {any $m$ between the two members of a pair costs exactly the pair's gap, and the median lies inside every pair};
\end{tikzpicture}

\caption{The pairing picture. Notice the constant: within a pair, the
two distances to $m$ always sum to the pair's gap. The median sits inside
\emph{every} pair, so it inherits the minimum possible total.}
\label{fig:median}
\end{figure}

\begin{example}[CSES \emph{Apartments} --- greedy stays ahead]
\label{ex:apts}
Applicant $i$ accepts $[a_i - k, a_i + k]$. Sort demands and sizes, walk
two pointers: skip apartments too small for the current applicant; skip
applicants the current apartment cannot satisfy; on a fit, match.
\emph{Stays-ahead invariant:} processing applicants in increasing demand,
after greedy's $i$-th step it has matched at least as many applicants as
any optimal solution restricted to the same first $i$ --- so when greedy
stops, no unseen match can exist anywhere in the input.
\lstinputlisting{code/apartments.cpp}
Real output on the CSES sample (\texttt{4 3 5 / 60 45 80 60 / 30 60 75}):
\begin{lstlisting}[numbers=none]
2
\end{lstlisting}
\end{example}

\begin{exercise}[easy]
On input \texttt{n=3, m=3, k=0}, demands $\{5,5,9\}$, sizes
$\{5,9,9\}$ --- does sorting the apartments change the greedy answer?
Answer first, then check by hand.
\end{exercise}

\section{When greedy fails}
\label{sec:fail}
\index{greedy algorithm!failure}

\begin{example}[Coin systems --- prove or distrust]
\label{ex:coins}
Coin change, minimum coins for amount $A$ over a fixed system. For
$\{1,5,10,25\}$, ``always take the largest possible'' is optimal (a
canonical system --- provable). For $\{1,3,4\}$ the same rule takes
$4+1+1$ for $6$: three coins where $3+3$ suffices.
\emph{Strategy.} Never trust the rule because it survives samples ---
test the \emph{smallest} candidate failure, and if you cannot prove the
rule, run DP (Chapter~11) instead.
\end{example}
\lstinputlisting{code/coins.cpp}
Real output --- greedy vs the true optimum on both systems:
\begin{lstlisting}[numbers=none]
amount=6 coins={1,3,4} greedy=3 optimal=2
amount=99 coins={1,5,10,25} greedy=9 optimal=9
\end{lstlisting}
Figure~\ref{fig:cointree} shows where the wrong branch starts.

\begin{figure}[ht]\centering
% fig14-cointree.tex -- Figure 3.5: coin tree for amount 6, system {1,3,4}. Greedy takes 4 first.
% fstyle.tex -- shared TikZ styles for the figures of this book (palette from lifeos.sty).
\tikzset{%
  cell/.style={draw=ink, minimum width=7mm, minimum height=6.5mm, inner sep=0pt, font=\footnotesize, fill=white},
  cellh/.style={cell, fill=accentsoft, draw=accent},
  cella/.style={cell, fill=amberbg, draw=amber},
  cellg/.style={cell, fill=goodbg, draw=good},
  cellx/.style={cell, fill=graybg, draw=slate, text=slate},
  vtx/.style={circle, draw=ink, minimum size=6.5mm, inner sep=0pt, font=\footnotesize, fill=white},
  vtxh/.style={vtx, draw=accent, fill=accentsoft},
  vtxa/.style={vtx, draw=amber, fill=amberbg},
  vtxg/.style={vtx, draw=good, fill=goodbg},
  vtxd/.style={vtx, draw=slate, fill=graybg, text=slate},
  lab/.style={font=\footnotesize, text=slate},
  labA/.style={font=\footnotesize, text=accent},
  labM/.style={font=\footnotesize, text=amber},
  idx/.style={font=\tiny, text=slate},
  arr/.style={-{stealth}, thick, ink},
  arrA/.style={-{stealth}, thick, accent},
  arrM/.style={-{stealth}, thick, amber},
  arrG/.style={-{stealth}, thick, good},
  edge/.style={thick, ink},
  edgeA/.style={thick, accent},
  edgeM/.style={thick, amber},
}%
% \cellrow{prefix}{x0}{y}{style}{v0,v1,...}: a row of cells named prefix0, prefix1, ... (pitch 0.75cm)
\providecommand{\cellrow}[5]{\foreach \v [count=\i from 0] in {#5} {\node[#4] (#1\i) at (#2+\i*0.75,#3) {\v};}}%
% \cellidx{prefix}{n}{start}: index labels under cells prefix0 .. prefix(n-1)
\providecommand{\cellidx}[3]{\foreach \i in {0,...,\numexpr#2-1\relax} {\node[idx, below=0pt of #1\i] {\the\numexpr\i+#3\relax};}}%
% \fitfig{<tikzpicture>}: shrink the picture to \linewidth only when it would overflow.
\ifdefined\fitbox\else\newsavebox{\fitbox}\fi%
\providecommand{\fitfig}[1]{\sbox{\fitbox}{#1}\ifdim\wd\fitbox>\linewidth\resizebox{\linewidth}{!}{\usebox{\fitbox}}\else\usebox{\fitbox}\fi}%

\begin{tikzpicture}[font=\footnotesize, x=1cm, y=1cm,
  el/.style={font=\scriptsize, text=slate, fill=white, inner sep=1pt}]
\node[vtx] (r) at (0,0) {6};
\node[vtxa] (a) at (-3.2,-1.5) {2};
\node[vtx]  (b) at (0,-1.5) {5};
\node[vtxg] (c) at (3.2,-1.5) {3};
\node[vtxa] (a2) at (-3.2,-3.0) {1};
\node[vtxa] (a3) at (-3.2,-4.5) {0};
\node[vtxg] (c2) at (3.2,-3.0) {0};
\node[lab] (bd) at (0,-3.0) {$\vdots$};
\draw[edgeM] (r) -- (a) node[el, midway, above left] {$-4$};
\draw[edge]  (r) -- (b) node[el, midway, right] {$-1$};
\draw[edgeA, good] (r) -- (c) node[el, midway, above right] {$-3$};
\draw[edgeM] (a) -- (a2) node[el, midway, left] {$-1$};
\draw[edgeM] (a2) -- (a3) node[el, midway, left] {$-1$};
\draw[edge, dashed, slate] (b) -- (bd);
\draw[edgeA, good] (c) -- (c2) node[el, midway, right] {$-3$};
\node[labM, anchor=west, align=left] at (-2.9,-4.5) {greedy: $4+1+1 = 3$ coins};
\node[font=\footnotesize, text=good, anchor=west, align=left] at (3.6,-3.0) {optimal: $3+3 = 2$ coins};
\end{tikzpicture}

\caption{The decision tree for amount $6$ on $\{1,3,4\}$. Notice the
failure is \emph{local}: the first greedy pick ($4$) already forfeits the
optimum --- there is no way back later. That is what ``greedy fails''
means mechanically.}
\label{fig:cointree}
\end{figure}

\begin{principle}[Prove or distrust]
\index{prove or distrust}
A greedy rule you cannot defend in two sentences gets a counter-example
search before it gets submitted. Try the smallest meaningful instance;
if it survives, try to \emph{prove} the rule survives; only then code it.
\end{principle}

\begin{exercise}[medium]
Find the smallest amount on which greedy fails for system
$\{1, 7, 10\}$. (Write the three-candidate check before guessing.)
\end{exercise}

\section{Sets and maps as greedy tools}
\label{sec:settools}
\index{std::set@\texttt{std::set}}

A greedy step often asks ``what is the best available option \emph{right
now}?'' --- which is exactly an ordered-set query:
\texttt{lower\_bound}/\texttt{upper\_bound} in $O(\log n)$, then
\texttt{erase} to consume the resource. CSES \emph{Concert Tickets} is
the archetype: sell each customer the most expensive seat $\le$ their
limit. Duplicates are real, so \texttt{multiset} --- and
\texttt{erase(it)} on an \emph{iterator} removes one copy, while
\texttt{erase(value)} removes them all.

\begin{lstlisting}[label={lst:setpat},
caption={\texttt{code/set\_pattern.cpp}: the ``best available'' idiom.
\texttt{upper\_bound(t)} finds the first seat $> t$; stepping back one
gives the largest $\le t$.}]
\end{lstlisting}
\lstinputlisting{code/set_pattern.cpp}
Real output on \texttt{5 3 / 5 3 7 8 5 / 4 8 3} --- seat $3$ was already
sold when the last customer asked:
\begin{lstlisting}[numbers=none]
3
8
-1
\end{lstlisting}

\section{Chapter summary}
\begin{table}[ht]\centering
{\small\begin{tabularx}{\linewidth}{@{}lYYY@{}}
\toprule
\thead{Situation} & \thead{Greedy rule} & \thead{Proof tool} & \thead{When to distrust}\\
\midrule
intervals/schedules & earliest \emph{end} & exchange argument & earliest start fails\\
equalise to a point & take the median & pairing (constant per pair) & weighted targets\\
two sorted lists & advance whichever pointer is behind & stays-ahead invariant & unsorted input\\
``best available'' & \texttt{set::lower\_bound} & step dominance & duplicates need \texttt{multiset}\\
coin-like systems & largest first & only if canonical & any non-canonical set\\
\bottomrule
\end{tabularx}}
\caption{The chapter in one table: the rule, the proof that licenses it,
and the smell that says stop. The middle column is the one to memorise.}
\label{tab:greedy}
\end{table}

\section{Common mistakes}
\index{common mistakes}
\begin{itemize}
\item \textbf{Comparator that is not a strict weak order.} \texttt{<=},
or a comparator that sometimes says both orders, corrupts
\texttt{sort}'s internal invariant --- silent wrong answers \emph{and}
possible crashes.
\item \textbf{Sorting the wrong key.} Movie Festival by start time gives
$1$ on a 2-answer case. The sort key \emph{is} the greedy rule --- pick
it by the proof, not by habit.
\item \textbf{Trusting greedy because samples pass.} Two samples cannot
witness a coin system's failure at amount $6$. Counter-example hunt
first, or proof, or DP.
\item \textbf{Forgetting stability.} \texttt{sort} reorders equals;
\texttt{stable\_sort} does not. Only matters when ties carry meaning ---
exactly when you'll least expect it.
\item \textbf{\texttt{erase(value)} on a multiset.} Removes \emph{every}
copy. Erase through the iterator the lookup returned.
\item \textbf{Compressing before seeing all values.} Ranks are defined
by the full set; an online stream needs a balanced structure, not
compression.
\end{itemize}

\section{Additional exercises}
\begin{exercise}[easy] CSES \emph{Distinct Numbers} --- what does sorting
buy you here that a \texttt{set} also gives, and which is faster to
write?
\end{exercise}
\begin{exercise}[easy] CSES \emph{Sum of Two Values} --- after sorting,
which two pointers, and when do they move? (This is the apartment
pattern without the tolerance.)
\end{exercise}
\begin{exercise}[easy] CSES \emph{Collecting Numbers} --- sort values
keeping original indices; count how many rounds a left-to-right pass
needs. Which adjacency test detects a new round?
\end{exercise}
\begin{exercise}[easy] USACO-B \emph{Cow College} --- the optimal tuition
is one of the offered prices; why does sorting make trying them all
$O(n\log n)$ instead of $O(n^2)$?
\end{exercise}
\begin{exercise}[easy] USACO-B \emph{Even More Odd Photos} --- greedy
rules and parity counting meet; identify the greedy rule and the one
case it must handle carefully.
\end{exercise}
\begin{exercise}[easy] USACO-S \emph{Lemonade Line} --- cows tolerate a
queue position only if enough room remains; sort and simulate. Prove the
greedy order is optimal by exchange.
\end{exercise}
\begin{exercise}[medium] CSES \emph{Restaurant Customers} --- sweep line
is sorting in disguise: sort all events together and walk a counter.
Implement it; then say why ``maximum prefix sum of the event order'' is
the answer.
\end{exercise}
\begin{exercise}[medium] CSES \emph{Tasks and Deadlines} --- which order
minimises total finishing time weighted equally? (Shortest processing
time first --- prove it by exchanging one adjacent inversion.)
\end{exercise}
\begin{exercise}[medium] USACO-S \emph{High Card Wins} --- Bessie's
strongest card should beat Elsie's weakest beatable card. Write the
stays-ahead invariant for this greedy.
\end{exercise}
\begin{exercise}[medium] USACO-S \emph{Rest Stops} --- Farmer John stops
only at tastiest-ahead stops. Sort or scan for ``record'' tastiness;
which direction does the scan run?
\end{exercise}
\begin{exercise}[medium] USACO-B \emph{Out of Place} --- one cow is
misplaced; sorting reveals her. Prove the answer is the count of cows
off their sorted position, minus one boundary case --- find that case.
\end{exercise}
\begin{exercise}[hard] CSES \emph{Josephus Problem I} --- simulate the
elimination circle. Greedy helps nothing; the honest answer is a
\texttt{set}/list walk or an order-statistics structure --- implement one
and compare with $O(n^2)$ simulation on $n=2\cdot10^5$.
\end{exercise}

\chapternotes
The greedy-proof taxonomy (exchange, stays-ahead, induction) and the
canonical/non-canonical coin example follow
Sannemo~\cite{sannemo2025} and Laaksonen's chapter~6~\cite{laaksonen2024};
the interval-scheduling exchange is the standard one, retold with the
counter-example-first pedagogy of the USACO Guide~\cite{usacoguide}.
\texttt{std::sort}'s introsort and strict-weak-ordering contract are
from cppreference~\cite{cppreference}; Halim's treatment of sort-and-greedy
problem patterns informed the exercise set~\cite{halimcp4}. All problem
names and samples are from the CSES Problem Set~\cite{cses}.

\section*{Further reading}
\begin{itemize}
\item Laaksonen, \emph{Competitive Programmer's Handbook}, ch.~6 ``Greedy
algorithms'' (beginner). \url{https://cses.fi/book/book.pdf} --- the
scheduling/median/coin-canonical results quoted in this chapter.
\item USACO Guide, ``Introduction to Greedy Algorithms'' (beginner).
\url{https://usaco.guide/bronze/intro-greedy} --- more greedy exercises at
the Bronze level, same decision discipline.
\item USACO Guide, ``Introduction to Sorting'' (beginner).
\url{https://usaco.guide/bronze/intro-sorting} --- comparators and
sorting-based problem solving beneath \S\ref{sec:sorting}.
\item CSES 1629 \emph{Movie Festival} (practice).
\url{https://cses.fi/problemset/task/1629} --- submit the earliest-finish
schedule of \ref{ex:movie} yourself.
\item cppreference, ``std::sort'' (reference).
\url{https://en.cppreference.com/w/cpp/algorithm/sort} --- comparator
contract (strict weak ordering) behind \S\ref{sec:sorting}'s rules.
\item Wikipedia, ``Greedy algorithm'' (reference).
\url{https://en.wikipedia.org/wiki/Greedy_algorithm} --- general context
for the exchange/stays-ahead proofs of \S\ref{sec:proofs}.
\end{itemize}

\begin{rulebox}[Watch instead]
\smallskip\noindent\textbf{Best first.}
``3. Greedy Method - Introduction'' --- Abdul Bari, 12 min ---
\url{https://www.youtube.com/watch?v=ARvQcqJ_-NY}. The greedy-choice
framing of \S\ref{sec:greedy} with small worked cases; he optimizes
objectives on the same ``choice + feasibility'' axis this book uses.
\smallskip\noindent
``Lec 16'' --- MIT OpenCourseWare 6.046J, 84 min ---
\url{https://www.youtube.com/watch?v=FPEMBWg_WlY}. Formal greedy-proof
machinery (greedy-choice property + induction) matching
\S\ref{sec:proofs}; blackboard pace, uses more algorithm design notation
than the book needs. Access date 2026-10-10.
\end{rulebox}
